\documentclass[11pt]{article}

\usepackage[margin=1in]{geometry}
\usepackage{amsmath,amssymb,amsthm}
\usepackage[hidelinks]{hyperref}

\newtheorem{theorem}{Theorem}[section]
\newtheorem{lemma}[theorem]{Lemma}
\newtheorem{proposition}[theorem]{Proposition}
\newtheorem{corollary}[theorem]{Corollary}
\newtheorem{conjecture}[theorem]{Conjecture}
\theoremstyle{definition}
\newtheorem{definition}[theorem]{Definition}
\theoremstyle{remark}
\newtheorem{remark}[theorem]{Remark}

\DeclareMathOperator{\ord}{ord}
\DeclareMathOperator{\lcm}{lcm}
\DeclareMathOperator{\dlog}{dlog}

\title{Partial Results on the Generalized Two-Dimensional Sierpi\'nski Problem}
\author{Animish Sharma}
\date{21 June 2026}

\begin{document}

\maketitle

\begin{abstract}
We study the first nontrivial two-prime case of Erd\H{o}s--Graham Problem
\#203: whether there exists an integer \(m\geq 1\), with \(\gcd(m,6)=1\),
such that
\[
        2^k 3^\ell m+1
\]
is composite for every pair \(k,\ell\geq 0\).  The central existence question
remains open.  This report records partial results toward the problem:
covering-set criteria, density necessary conditions for both linear and
nonlinear primes, an impossibility hierarchy for small covering moduli, the
structural implication that every two-dimensional Sierpi\'nski number is also a
classical one-dimensional Sierpi\'nski number, and computational evidence that
the density condition is not sufficient for a two-dimensional covering.  The
latest WIT version incorporated here is WIT v10, which extends the main
mixed-prime LCM obstruction to all primes \(p\leq 50{,}000{,}000\).
\end{abstract}

\section{Problem Statement and Status}

The problem considered here is:

\begin{quote}
Does there exist \(m\geq 1\), with \(\gcd(m,6)=1\), such that
\[
        2^k3^\ell m+1
\]
is composite for all \(k,\ell\geq 0\)?
\end{quote}

This is the \(r=2\), \(p_1=2\), \(p_2=3\) case of Erd\H{o}s--Graham Problem
\#203, which asks whether for given primes \(p_1,\ldots,p_r\) there is an
integer \(m\) such that
\[
        p_1^{k_1}\cdots p_r^{k_r}m+1
\]
is composite for all \(k_1,\ldots,k_r\geq 0\).

\medskip
\noindent
\textbf{Current status.}  The existence question remains open.  No proof of
existence and no proof of nonexistence is claimed in this report.

\section{Definitions}

\begin{definition}[Linear prime]
Let \(p\) be a prime with \(\gcd(p,6)=1\).  We call \(p\) \emph{linear} if
\[
        3\in \langle 2\rangle \subseteq (\mathbb Z/p\mathbb Z)^\times.
\]
Equivalently, there exists \(e\in\{0,\ldots,\ord_p(2)-1\}\) such that
\[
        2^e\equiv 3 \pmod p.
\]
In this case \(e=\dlog_2(3)\pmod p\) is defined relative to the cyclic subgroup
\(\langle 2\rangle\).
\end{definition}

\begin{definition}[Diagonal strip]
Let \(p\) be a linear prime, set \(a=\ord_p(2)\), and let
\(e=\dlog_2(3)\pmod p\).  For a residue \(r\in\{0,\ldots,a-1\}\), define
\[
        S(p,r)=\{(k,\ell)\in \mathbb Z^2: k+e\ell\equiv r \pmod a\}.
\]
If \(m\) is coprime to \(p\), then
\[
        p\mid 2^k3^\ell m+1
\]
if and only if
\[
        (k,\ell)\in S(p,r_m),
        \qquad
        r_m=\dlog_2(-m^{-1}\bmod p).
\]
\end{definition}

\begin{definition}[Subgroup attached to a prime]
For every prime \(p\) with \(\gcd(p,6)=1\), define
\[
        H_p=\langle 2,3\rangle\subseteq (\mathbb Z/p\mathbb Z)^\times.
\]
For fixed \(m\) coprime to \(p\), put \(c=-m^{-1}\pmod p\).  Then
\[
        p\mid 2^k3^\ell m+1
        \quad\Longleftrightarrow\quad
        2^k3^\ell\equiv c\pmod p.
\]
If \(c\notin H_p\), the prime \(p\) covers no pairs \((k,\ell)\) for that
\(m\).  If \(c\in H_p\), the covered set has density \(1/|H_p|\) on its
period grid.
\end{definition}

\begin{definition}[Two-dimensional covering set]
A finite set \(P\) of primes with \(\gcd(p,6)=1\), together with compatible
residue choices for each \(p\in P\), is called a \emph{two-dimensional covering
set for \(m\)} if for every \(k,\ell\geq 0\) some \(p\in P\) satisfies
\[
        p\mid 2^k3^\ell m+1.
\]
Such a covering set implies that \(m\) is a two-dimensional generalized
Sierpi\'nski number, provided each divisor is a proper divisor of
\(2^k3^\ell m+1\).
\end{definition}

\section{Basic Group-Theoretic Structure}

\begin{lemma}[Order of the generated subgroup]
Let \(p\) be a prime with \(\gcd(p,6)=1\).  Let
\[
        a=\ord_p(2),\qquad b=\ord_p(3).
\]
Then
\[
        |H_p|=|\langle 2,3\rangle|=\lcm(a,b).
\]
\end{lemma}

\begin{proof}
The group \((\mathbb Z/p\mathbb Z)^\times\) is cyclic of order \(p-1\).  Let
\(g\) be a primitive root modulo \(p\), and write
\[
        2\equiv g^\alpha,\qquad 3\equiv g^\beta \pmod p.
\]
Then
\[
        a=\frac{p-1}{\gcd(\alpha,p-1)},\qquad
        b=\frac{p-1}{\gcd(\beta,p-1)}.
\]
In the additive cyclic group \(\mathbb Z/(p-1)\mathbb Z\), the subgroup
generated by \(\alpha\) and \(\beta\) is generated by
\(\gcd(\alpha,\beta)\).  Hence
\[
        |\langle 2,3\rangle|
        =
        \frac{p-1}{\gcd(\alpha,\beta,p-1)}.
\]
On the other hand, using the identity
\[
        \lcm\left(\frac n{d_1},\frac n{d_2}\right)
        =
        \frac n{\gcd(d_1,d_2)}
        \qquad(d_1,d_2\mid n),
\]
with \(n=p-1\), \(d_1=\gcd(\alpha,p-1)\), and
\(d_2=\gcd(\beta,p-1)\), we get
\[
        \lcm(a,b)
        =
        \frac{p-1}{\gcd(\alpha,\beta,p-1)}.
\]
Thus \(|H_p|=\lcm(a,b)\).
\end{proof}

\section{Density Lemmas}

\begin{lemma}[Diagonal strip density]
Let \(a,L\geq 1\), with \(a\mid L\), and let \(e,r\in\{0,\ldots,a-1\}\).
Then the number of pairs \((k,\ell)\in\{0,\ldots,L-1\}^2\) satisfying
\[
        k+e\ell\equiv r\pmod a
\]
is exactly \(L^2/a\).
\end{lemma}

\begin{proof}
For each fixed \(\ell\), the congruence determines exactly one residue class
for \(k\) modulo \(a\), namely
\[
        k\equiv r-e\ell\pmod a.
\]
Since \(a\mid L\), the interval \(\{0,\ldots,L-1\}\) contains exactly \(L/a\)
representatives of that residue class.  Summing over all \(L\) choices of
\(\ell\), the total number of solutions is \(L(L/a)=L^2/a\).
\end{proof}

\begin{lemma}[Subgroup density]
Let \(p\) be a prime with \(\gcd(p,6)=1\), let
\[
        H_p=\langle 2,3\rangle,\qquad h=|H_p|,
\]
and let \(m\geq 1\) with \(\gcd(m,6)=1\).  Suppose
\[
        c=-m^{-1}\pmod p
\]
lies in \(H_p\).  If \(L\geq 1\) and \(h\mid L\), then the number of pairs
\((k,\ell)\in\{0,\ldots,L-1\}^2\) satisfying
\[
        p\mid 2^k3^\ell m+1
\]
is exactly \(L^2/h\).
\end{lemma}

\begin{proof}
Let \(a=\ord_p(2)\) and \(b=\ord_p(3)\).  The map
\[
        \phi:\mathbb Z/a\mathbb Z\times \mathbb Z/b\mathbb Z\to H_p,
        \qquad
        (k,\ell)\mapsto 2^k3^\ell
\]
is a well-defined surjective homomorphism.  Since \(|H_p|=h\), its kernel has
size \(ab/h\).  Therefore each element \(c\in H_p\) has exactly \(ab/h\)
preimages modulo \(a\) and \(b\).

Because \(h=\lcm(a,b)\), the condition \(h\mid L\) implies \(a\mid L\) and
\(b\mid L\).  Each solution modulo \(a,b\) lifts to exactly
\[
        (L/a)(L/b)
\]
pairs in the square \(\{0,\ldots,L-1\}^2\).  Thus the number of pairs is
\[
        \frac{ab}{h}\cdot \frac La\cdot\frac Lb=\frac{L^2}{h}.
\]
\end{proof}

\section{Covering Sets and Sierpi\'nski Numbers}

\begin{theorem}[Covering implies compositeness]
Let \(P\) be a finite set of primes with \(\gcd(p,6)=1\) for every
\(p\in P\), and let \(m\geq 1\) satisfy \(\gcd(m,6)=1\).  Suppose that for all
\(k,\ell\geq 0\), there exists \(p\in P\) such that
\[
        p\mid 2^k3^\ell m+1.
\]
If each such divisor is proper, then \(2^k3^\ell m+1\) is composite for all
\(k,\ell\geq 0\).
\end{theorem}

\begin{proof}
Fix \(k,\ell\geq 0\), and set
\[
        N=2^k3^\ell m+1.
\]
By hypothesis, some \(p\in P\) divides \(N\).  If \(1<p<N\), then \(N\) has a
proper prime divisor and is therefore composite.

In explicit CRT-derived covering constructions, the corresponding \(m\) is
chosen so that \(N>p\) for every prime in the covering set and every
\((k,\ell)\) in the finite period grid.  Equivalently, one verifies the finite
small-case condition \(N\neq p\) on the period grid.  This is automatic in the
large CRT constructions considered computationally below, where the resulting
\(m\) is much larger than \(\max P\).
\end{proof}

\begin{remark}
The finite ``proper divisor'' check is not a density issue.  It is a standard
CRT verification step: after a proposed residue assignment is found, one checks
on the finite period grid that the covering primes are proper factors of the
covered values.
\end{remark}

\section{A Lov\'asz Local Lemma Check}

\begin{lemma}[Failure of the symmetric Lov\'asz local lemma test]
Let \(P_{5040}\) be the 31-prime set whose associated subgroup sizes
\(|H_p|\) all divide \(5040\), and let \(N=5040\).  On the torus
\(\{0,\ldots,N-1\}^2\), choose one residue class independently for each prime.
For each cell \((k,\ell)\), let \(A_{k,\ell}\) be the event that no selected
prime-residue class covers \((k,\ell)\).  The symmetric Lov\'asz local lemma
condition
\[
        e\,\Pr(A_{k,\ell})(d+1)\leq 1
\]
fails by a large margin for this model.
\end{lemma}

\begin{proof}
For a fixed cell \((k,\ell)\), prime \(p_i\) covers the cell with probability
\(1/h_i\), where \(h_i=|H_{p_i}|\).  Hence
\[
        \Pr(A_{k,\ell})
        =
        \prod_i\left(1-\frac1{h_i}\right).
\]
The WIT computation records this probability as at least \(0.36\) for the
31-prime \(5040\)-set.

The dependency degree is also enormous.  Since prime \(p_i\) covers
\(N^2/h_i\) cells, a fixed noncoverage event can interact with on the order of
\[
        \sum_i \frac{N^2}{h_i}
        =
        N^2\sum_i \frac1{h_i}
\]
other events.  For the \(5040\)-set,
\[
        \sum_i \frac1{h_i}=1.021429\ldots,
\]
so the dependency degree is greater than \(25{,}000{,}000\).  Therefore
\[
        e\,\Pr(A_{k,\ell})(d+1)
        >
        2.7\cdot 0.36\cdot 25{,}000{,}000
        \gg 1.
\]
The symmetric local lemma therefore gives no existence guarantee for a full
covering assignment.
\end{proof}

\begin{remark}
This lemma is a negative result about a proof method only.  It does not prove
that the \(5040\)-set has no covering assignment.
\end{remark}

\section{Density Necessary Conditions}

\begin{theorem}[Density condition for linear primes]
Let \(P=\{p_1,\ldots,p_n\}\) be a finite set of linear primes, and set
\[
        a_i=\ord_{p_i}(2).
\]
If the corresponding diagonal strips cover all of \(\mathbb Z_{\geq 0}^2\),
then
\[
        \sum_{i=1}^n \frac1{a_i}\geq 1.
\]
\end{theorem}

\begin{proof}
Let \(L=\lcm(a_1,\ldots,a_n)\).  Each strip for \(p_i\) occupies exactly
\(L^2/a_i\) points in the square \(\{0,\ldots,L-1\}^2\), by the diagonal strip
density lemma.  Since the strips cover the whole square, the total number of
strip incidences is at least \(L^2\).  Therefore
\[
        \sum_{i=1}^n \frac{L^2}{a_i}\geq L^2.
\]
Dividing by \(L^2\) gives the result.
\end{proof}

\begin{theorem}[General density condition]
Let \(P=\{p_1,\ldots,p_n\}\) be a finite set of primes with \(\gcd(p_i,6)=1\).
Let
\[
        h_i=|H_{p_i}|=|\langle 2,3\rangle \bmod p_i|.
\]
Assume that for the fixed \(m\) under consideration the relevant target
\(-m^{-1}\pmod {p_i}\) lies in \(H_{p_i}\) for each \(i\), and that the
covering patterns
\[
        \{(k,\ell):p_i\mid 2^k3^\ell m+1\}
\]
cover all of \(\mathbb Z_{\geq 0}^2\).  Then
\[
        \sum_{i=1}^n \frac1{h_i}\geq 1.
\]
\end{theorem}

\begin{proof}
Let \(L=\lcm(h_1,\ldots,h_n)\).  By the subgroup density lemma, the pattern
corresponding to \(p_i\) covers exactly \(L^2/h_i\) points in
\(\{0,\ldots,L-1\}^2\).  Since the patterns cover the square,
\[
        \sum_{i=1}^n \frac{L^2}{h_i}\geq L^2.
\]
Dividing by \(L^2\) gives
\[
        \sum_{i=1}^n \frac1{h_i}\geq 1.
\]
\end{proof}

\begin{remark}
For linear primes, \(h_i=\ord_{p_i}(2)\), so the general theorem recovers the
linear density condition.  For nonlinear primes, the subgroup size is
\[
        h_i=\lcm(\ord_{p_i}(2),\ord_{p_i}(3)).
\]
\end{remark}

\section{Impossibility Results}

\begin{theorem}[The four-prime set fails]
No two-dimensional covering set exists using only the linear primes
\[
        \{5,11,23,13\}.
\]
\end{theorem}

\begin{proof}
The relevant orders are
\[
        \ord_5(2)=4,\quad
        \ord_{11}(2)=10,\quad
        \ord_{23}(2)=11,\quad
        \ord_{13}(2)=12.
\]
Thus
\[
        \frac14+\frac1{10}+\frac1{11}+\frac1{12}
        =
        \frac{165+66+60+55}{660}
        =
        \frac{346}{660}
        <1.
\]
This violates the density necessary condition for linear primes.
\end{proof}

\begin{theorem}[Linear primes of order at most \(60\) fail]
No two-dimensional covering set exists using only linear primes \(p\) with
\(\ord_p(2)\leq 60\).
\end{theorem}

\begin{proof}
The linear primes with \(\ord_p(2)\leq 60\) are recorded as
\[
        5,11,23,13,19,47,29,71,37,431,97,53,59,61,
\]
with orders
\[
        4,10,11,12,18,23,28,35,36,43,48,52,58,60.
\]
Their reciprocal sum is
\[
\begin{aligned}
&\frac14+\frac1{10}+\frac1{11}+\frac1{12}
+\frac1{18}+\frac1{23}+\frac1{28}+\frac1{35} \\
&\qquad
+\frac1{36}+\frac1{43}+\frac1{48}+\frac1{52}
+\frac1{58}+\frac1{60}
        =0.81256769\ldots<1.
\end{aligned}
\]
Every subset has still smaller reciprocal sum, so every such set violates the
density necessary condition.
\end{proof}

\begin{theorem}[Linear-prime LCM obstruction]
Among linear primes \(p<5000\), no covering set whose order LCM is at most
\[
        720720
\]
can exist.
\end{theorem}

\begin{proof}
For every \(L\leq 720720\), let
\[
        S_L=\{p<5000:p\text{ is linear and }\ord_p(2)\mid L\}.
\]
An exhaustive computation over \(L\leq 720720\) gives
\[
        \max_{L\leq 720720}\sum_{p\in S_L}\frac1{\ord_p(2)}
        =
        0.7930527805527805<1.
\]
Therefore every linear-prime set with LCM of orders at most \(720720\) violates
the density necessary condition.
\end{proof}

\begin{theorem}[Mixed-prime LCM obstruction below \(5040\), computational scope]
For prime sets with all primes \(p\leq 50{,}000{,}000\), no two-dimensional covering set
can have
\[
        \lcm\{|H_p|:p\in P\}<5040.
\]
\end{theorem}

\begin{proof}
For \(B<5040\), define
\[
        \sigma_{\mathrm{gen}}(B)
        =
        \sum_{\substack{p\leq 50{,}000{,}000\\ \gcd(p,6)=1\\ |H_p|\mid B}}
        \frac1{|H_p|}.
\]
The computation recorded in the WIT file exhaustively sieves all primes
\(p\leq 50{,}000{,}000\).  It finds that
\[
        \sigma_{\mathrm{gen}}(B)<1
        \qquad\text{for every }B<5040.
\]
The maximum below \(5040\) occurs at \(B=2520\), with value \(0.917460\).  The
Phase 6 extension found 10 new primes in the interval
\((10^6,50{,}000{,}000]\) with \(|H_p|\mid 5040\), but their contribution is
too small to make any \(\sigma_{\mathrm{gen}}(B)\) with \(B<5040\) reach \(1\).

If \(P\) had \(\lcm(|H_p|:p\in P)=B<5040\), then
\[
        \sum_{p\in P}\frac1{|H_p|}
        \leq \sigma_{\mathrm{gen}}(B)<1,
\]
contradicting the general density necessary condition.
\end{proof}

\begin{remark}
This theorem has a computational scope: it is proved for primes
\(p\leq 50{,}000{,}000\).  Primes \(p>50{,}000{,}000\) with
\(|H_p|\mid 5040\) would be large prime factors of \(2^{5040}-1\) and
\(3^{5040}-1\); the WIT notes that a full universal closure would require
factoring a very large integer or a new theoretical argument.
\end{remark}

\section{The Structural Reduction to the Classical Problem}

\begin{theorem}[Two-dimensional implies one-dimensional]
If \(m\geq 1\), \(\gcd(m,6)=1\), and
\[
        2^k3^\ell m+1
\]
is composite for all \(k,\ell\geq 0\), then
\[
        2^k m+1
\]
is composite for all \(k\geq 0\).
\end{theorem}

\begin{proof}
Set \(\ell=0\).  Then
\[
        2^k3^0m+1=2^km+1.
\]
Thus the two-dimensional property immediately implies the classical
one-dimensional Sierpi\'nski property.
\end{proof}

\begin{corollary}
Every two-dimensional generalized Sierpi\'nski number is a classical
Sierpi\'nski number.
\end{corollary}

\section{The Candidate \(m=78557\)}

The integer \(78557\) is the classical smallest known Sierpi\'nski number.  It
does not solve the two-dimensional problem.

\begin{theorem}
The integer \(m=78557\) is not a two-dimensional generalized Sierpi\'nski
number.
\end{theorem}

\begin{proof}
At \((k,\ell)=(1,2)\),
\[
        2^1 3^2\cdot 78557+1
        =
        18\cdot 78557+1
        =
        1414027.
\]
The number \(1414027\) is prime, by the computational primality check recorded
in the WIT evidence.  Hence \(2^k3^\ell m+1\) is not always composite for
\(m=78557\).
\end{proof}

\begin{remark}
Also \(78557=17\cdot 4621\), so \(p=17\) never divides
\[
        2^k3^\ell\cdot 78557+1,
\]
because the expression is always congruent to \(1\pmod {17}\).
\end{remark}

\section{Computational Evidence}

The following results are computational and should not be confused with a
formal resolution of the open problem.

\subsection{The 13-prime mixed set}

Including nonlinear primes, the greedy density threshold is reached by the
13-prime set
\[
        \{5,7,11,23,13,17,19,47,29,31,71,37,73\}.
\]
Its density is approximately
\[
        \sum_p\frac1{|H_p|}=1.0056,
\]
and
\[
        \lcm(|H_p|)=1275120.
\]
However, the true two-dimensional verification domain is not
\(\lcm(|H_p|)^2\).  The recorded periods are
\[
        A=\lcm(\ord_p(2))=637560,\qquad
        B=\lcm(\ord_p(3))=1275120,
\]
so
\[
        A B = 812965507200 \approx 8.13\cdot 10^{11}.
\]
This makes exhaustive SAT/SMT verification currently infeasible.  Thus the
13-prime set satisfies the density necessary condition, but its tileability
remains unresolved.

\subsection{The 31-prime \(5040\)-set}

A tractable set with \(A=B=5040\) is
\[
\begin{aligned}
P_{5040}=\{&
5,7,11,13,17,19,29,31,37,41,43,61,71,73,97,113,127,181,\\
&211,241,281,337,421,577,631,673,1009,2017,2521,3361,13441\}.
\end{aligned}
\]
For this set,
\[
        \sum_{p\in P_{5040}}\frac1{|H_p|}=1.021429>1,
\]
and the fundamental domain has
\[
        5040^2=25401600
\]
cells.

Greedy set cover with 30 random restarts gave maximum coverage
\[
        18216600/25401600=71.714\%.
\]
The number of uncovered cells was
\[
        7185000,
\]
which is exactly
\[
        \frac{7185000}{25401600}=\frac{99}{350}.
\]
No full covering assignment was found.

\begin{remark}
This is strong evidence that the density condition is not sufficient for
two-dimensional covering.  Strictly speaking, greedy failure alone is not a
formal proof that no assignment exists, but later ILP and simulated annealing
attempts also failed to find a covering.
\end{remark}

\subsection{ILP, simulated annealing, and Lov\'asz local lemma checks}

For the \(5040\)-set, an ILP formulation used binary variables
\[
        y_{p,r}\in\{0,1\},
\]
where \(y_{p,r}=1\) means that prime \(p\) uses residue \(r\).  The recorded
model had \(9279\) binary variables and thousands of seed covering constraints.
CBC did not find an integer feasible solution in the recorded runs, but did
not produce a formal infeasibility certificate.

Simulated annealing with \(1.5\cdot 10^6\) iterations reached
\[
        71.570\%
\]
coverage, below the greedy maximum \(71.714\%\).

The formal WIT v10 lemma \emph{lovasz-lll-fails} records the corresponding
cell-event local lemma calculation.  For a fixed cell \((k,\ell)\), the
noncoverage event has probability
\[
        \prod_i\left(1-\frac1{h_i}\right),
\]
and the dependency degree is on the order of
\[
        5040^2\sum_i\frac1{h_i}>25{,}000{,}000.
\]
Consequently the symmetric LLL condition fails by many orders of magnitude.

\subsection{The \(10080\)-modulus search}

The next tractable modulus considered was
\[
        N=10080=2\cdot 5040.
\]
The corresponding search used 33 primes, with density
\[
        1.031944.
\]
The best greedy coverage was
\[
        72.07\%=
        73223308/101606400.
\]
No covering was found.  The marginal improvement from \(5040\) to \(10080\) was
only about \(0.36\) percentage points.

\subsection{Phase 6: extension to \(p\leq 50{,}000{,}000\)}

The Phase 6 computation extended the T6b sieve from \(p\leq 10^6\) to
\[
        p\leq 50{,}000{,}000.
\]
It scanned \(3{,}932{,}436\) primes in \((10^6,50{,}000{,}000]\), testing
\[
        2^{5040}\equiv 1\pmod p,\qquad
        3^{5040}\equiv 1\pmod p.
\]
It found 10 new primes with \(|H_p|\mid 5040\):
\[
\begin{gathered}
2035153,\ 2543689,\ 3391249,\ 5084857,\ 6779137,\\
8473081,\ 13561969,\ 27110497,\ 33884761,\ 45175201.
\end{gathered}
\]
Their subgroup sizes are
\[
        1008,\ 2520,\ 5040,\ 5040,\ 5040,\ 5040,\ 5040,\ 5040,\ 5040,\ 1680,
\]
respectively.  The total density increment is
\[
        0.003373.
\]
Thus
\[
        \sigma_{\mathrm{gen}}(5040)=1.024802
\]
for the 41 primes found up to \(50{,}000{,}000\), while
\[
        \sigma_{\mathrm{gen}}(B)<1
        \qquad(1\leq B<5040).
\]
The maximum below \(5040\) remains at \(B=2520\), now with value
\[
        \sigma_{\mathrm{gen}}(2520)=0.917460.
\]

The same Phase 6 run also performed a greedy and local-search covering attempt
on a 23-prime subset of density \(0.944\).  The best coverage was
\[
        17683759/25401600=69.62\%,
\]
and no covering was found.

\section{Summary of Main Contributions}

\begin{enumerate}
    \item The problem is identified as the \(r=2\), \((p_1,p_2)=(2,3)\) case
    of Erd\H{o}s--Graham Problem \#203.

    \item A general density necessary condition is proved:
    every finite two-dimensional covering set must satisfy
    \[
        \sum_{p\in P}\frac1{|H_p|}\geq 1,
        \qquad
        H_p=\langle 2,3\rangle\bmod p.
    \]

    \item The subgroup size is proved generally:
    \[
        |H_p|=\lcm(\ord_p(2),\ord_p(3)).
    \]

    \item Several families of covering sets are ruled out by density:
    the four-prime set \(\{5,11,23,13\}\), all linear primes with
    \(\ord_p(2)\leq 60\), linear-prime sets with LCM at most \(720720\), and
    mixed-prime sets with \(p\leq 50{,}000{,}000\) and
    \(\lcm(|H_p|)<5040\).

    \item Any two-dimensional generalized Sierpi\'nski number is necessarily a
    classical one-dimensional Sierpi\'nski number.

    \item The natural candidate \(m=78557\) is ruled out by the prime value
    \[
        2^1 3^2\cdot 78557+1=1414027.
    \]

    \item Computational evidence shows that density greater than \(1\) is not
    by itself enough to produce a two-dimensional covering: the 31-prime
    \(5040\)-set has density \(1.021429\), but extensive search reached only
    about \(71.7\%\) coverage.

    \item Phase 6 extends the \(B^*=5040\) computation to all primes
    \(p\leq 50{,}000{,}000\): \(\sigma_{\mathrm{gen}}(B)<1\) for all
    \(B<5040\), while \(\sigma_{\mathrm{gen}}(5040)=1.024802\).
\end{enumerate}

\section{Version History of the WIT Development}

\begin{center}
\begin{tabular}{llll}
\hline
Version & Commit & Steps/Lines & Main addition \\
\hline
v1 & 19d70c6 & 20 / 242 & Linear-prime framework \\
v2 & c04460f & 31 / 455 & General \(H_p\) theorem \\
v3 & db4a6b3 & 34 / 596 & \(5040\)-set SAT evidence \\
v4 & 05d0ed3 & 34 / 648 & Uncovered-structure analysis \\
v5 & 22598e6 & 38 / 702 & Cyclic-subgroup proof and T1 fix \\
v6 & 3fbf6ba & 42 / 940 & T6b and GAP4 evidence \\
v6-final & de7ca08 & 42 / 1005 & \(N=10080\) evidence \\
v7 & f482a25 & 42 / 1069 & Scope and Lov\'asz documentation \\
v8 & caacb14 & 42 / \(\sim1120\) & T6b extended to \(p\leq 10^6\) \\
v9 & b11f534 & 45 / \(\sim1260\) & Formal Lov\'asz lemma and verification note \\
v10 & ec217a3 & 45 / \(\sim1340\) & T6b extended to \(p\leq 50{,}000{,}000\) \\
\hline
\end{tabular}
\end{center}

\section{WIT v10 Status}

The current WIT artifact is recorded as WIT v10, commit \texttt{ec217a3}.
It contains 45 top-level WIT steps and approximately 1340 lines.  Its structural
contents are:

\begin{center}
\begin{tabular}{ll}
\hline
Block type & Contents \\
\hline
Definitions & linear-prime, diagonal-strip, 2d-covering-set, subgroup-prime \\
Lemmas & cyclic-subgroup-order, diagonal-strip-density, subgroup-density, lovasz-lll-fails \\
Theorems & covering, density, impossibility, T6b, 2D-to-1D, \(m=78557\) \\
Conjecture & two-dimensional Sierpi\'nski existence, marked STATUS OPEN \\
Evidence & \(5040\), \(10080\), T6b universal, and Phase 6 computations \\
\hline
\end{tabular}
\end{center}

The WIT checker output recorded in the file is:
\[
        \text{45 top-level steps, 0 sub-steps, all references valid.}
\]
The same file also includes a human-verifiable checklist for readers who do not
have the WIT checker available.  A fresh critic session for WIT v10 reported
``complete, ship'' after checking the five requested items: open status of the
main conjecture, T6b scope \(p\leq 50{,}000{,}000\), the Phase 6 evidence
block, the formal Lov\'asz lemma, and the citation of the small-case remark in
the covering theorem.

\section{Open Problems}

\begin{conjecture}[Two-dimensional generalized Sierpi\'nski existence]
There exists \(m\geq 1\), with \(\gcd(m,6)=1\), such that
\[
        2^k3^\ell m+1
\]
is composite for all \(k,\ell\geq 0\).
\end{conjecture}

This conjecture remains open.  The present work provides obstructions and
computational evidence, but neither constructs such an \(m\) nor proves that no
such \(m\) can exist.

\section*{References}

\begin{enumerate}
    \item A. Granville and F. Pappalardi, \emph{Two Dimensional Covering
    Systems}, arXiv:2601.10296, 2026.

    \item J. Cremona and P. Koymans, \emph{Lattice Coverings and Homogeneous
    Covering Congruences}, arXiv:2601.03212, 2026.

    \item K. Hashimoto, \emph{Finiteness of Disjoint Covering Systems},
    arXiv:2603.26043, 2026.

    \item P. Erd\H{o}s, \emph{On integers of the form \(2^k+p\) and related
    problems}, 1950.

    \item P. Erd\H{o}s and R. L. Graham, \emph{Old and New Problems and Results
    in Combinatorial Number Theory}, Monographie No. 28 de L'Enseignement
    Math\'ematique, 1980.
\end{enumerate}

\end{document}
